\documentclass[11pt]{article}
\usepackage{explorer}
\newcommand{\ExplorerLanguage}{EN}
\newcommand{\ExplorerNumber}{01}
\newcommand{\ExplorerKind}{Lesson and method}
\begin{document}
\ExplorerTitle{High school · reasoning}{Parity and invariants}
\ExplorerSection{Before you begin}
Even and odd integers; proof by contradiction.\par
\ExplorerSection{Ideas and method}
An invariant is a quantity or property preserved by every allowed move. If the target has a different invariant value from the initial state, it cannot be reached.\par
To study parity, replace integers with their remainders 0 or 1 modulo 2. Two odd integers have an even sum. Adding an even integer preserves parity.\par
Method: specify one move, find what it preserves, and compute the invariant at the start and at the target. An obstruction proves impossibility; the absence of an obstruction does not prove possibility.\par
\ExplorerSection{Worked example}
Start at 0 and add or subtract 6 at each move. Every reachable number is even, so 15 is unreachable. The number 14 is also unreachable, but parity alone cannot prove this: every reachable number is a multiple of 6.\par
\ExplorerSection{Explore}
Write several sequences of moves, then classify the reached numbers. Find a common property more precise than parity.\par
\end{document}
